\documentclass[12pt,a4paper]{article}

\usepackage[utf8]{inputenc}
\usepackage[T1]{fontenc}
\usepackage{amsmath,amssymb}
\usepackage{booktabs}
\usepackage{hyperref}
\usepackage[margin=1in]{geometry}
\usepackage{natbib}
\usepackage{xcolor}
\usepackage{array}
\usepackage{caption}

\hypersetup{
    colorlinks=true,
    linkcolor=blue!60!black,
    citecolor=green!50!black,
    urlcolor=blue!70!black,
}

\title{Decentralized Cooperative Positioning via Trust-Weighted Fusion:\\
A Peer-to-Peer Alternative to GNSS on Commodity Hardware}

\author{
Tristan Stoltz\textsuperscript{1}\\[4pt]
\textsuperscript{1}Luminous Dynamics\\
\texttt{tristan.stoltz@evolvingresonantcocreationism.com}
}

\date{March 2026}

\begin{document}
\maketitle

% ============================================================================
% Abstract
% ============================================================================

\begin{abstract}
We present a decentralized positioning system that maintains location awareness when GNSS is denied, using only commodity smartphone sensors (BLE, WiFi RTT, UWB, accelerometer, gyroscope, barometer) and peer-to-peer ranging between cooperating devices. The system requires no dedicated infrastructure: phones with GNSS act as ephemeral anchors for phones without, and when all satellite signals are lost, Pedestrian Dead Reckoning (PDR) from the phone's inertial measurement unit maintains positioning for 5--30 minutes depending on magnetometer availability.

The architecture consists of three layers: (1)~a pure-Rust positioning library implementing Gauss-Newton weighted trilateration, an Extended Kalman Filter with barometric altitude, and adaptive-stride PDR (Weinberg model); (2)~a Holochain-based decentralized anchor registry with consciousness-gated Sybil resistance; and (3)~a mobile integration bridge supporting 8 ranging technologies including BLE Channel Sounding (Bluetooth 6.0, centimeter accuracy).

Benchmarks on release-compiled Rust show 3D trilateration at 11.5~\textmu s/call, PDR at 0.3~\textmu s/step (3.3M~steps/sec), and EKF convergence to sub-meter accuracy in 80 updates. The system is body-agnostic: the same algorithms work on Earth (WGS-84), Moon (IAU sphere), and Mars (IAU ellipsoid), making it applicable to lunar colonies and generation ships where no GNSS constellation exists.

The key contribution is the trust model: each range measurement is weighted by the reporting node's reputation score, derived from a four-dimensional profile (identity, reputation, community, engagement) stored on a Holochain distributed hash table. This provides Sybil resistance without a central authority---a property absent from all existing cooperative positioning proposals.
\end{abstract}

\textbf{Keywords:} decentralized positioning, cooperative localization, GNSS-denied navigation, trust-weighted fusion, Sybil resistance, pedestrian dead reckoning, Holochain, BLE Channel Sounding

% ============================================================================
% 1. Introduction
% ============================================================================

\section{Introduction}

Global Navigation Satellite Systems (GNSS) are a single point of failure embedded in nearly every layer of modern infrastructure. In May 2024, a solar storm disrupted GPS-based tractor guidance systems across the American Midwest, halting spring planting overnight. The U.S. Department of Homeland Security estimates that a nationwide GPS outage would cost the economy over \$1 billion per day \citep{dhs2021pnt}. A Carrington-class coronal mass ejection---an event with estimated recurrence of 100--200 years---would degrade or destroy the entire GNSS constellation for hours to days \citep{cannon2013extreme}.

Existing alternatives to GNSS (eLoran, Iridium STL, NextNav) share a critical limitation: they require centralized infrastructure controlled by a single entity. When that entity fails, the positioning service fails with it.

We propose and implement a fundamentally different approach: \emph{decentralized cooperative positioning}, where participating devices collectively establish position through peer-to-peer ranging, trust-weighted measurement fusion, and distributed anchor management. No central server, no dedicated infrastructure, no single point of failure.

The system addresses three scenarios:
\begin{enumerate}
    \item \textbf{GPS augmentation}: When GNSS is available, phones with GPS act as anchors for phones without (indoor users, underground, urban canyons).
    \item \textbf{GPS denial}: When GNSS is lost (solar storm, jamming, spoofing), the mesh of recently-positioned devices continues providing positioning with graceful degradation.
    \item \textbf{No-GNSS environments}: Lunar colonies, Mars settlements, and generation ships where no satellite constellation exists---peer-to-peer positioning from day one.
\end{enumerate}

\subsection{Contributions}

\begin{enumerate}
    \item A pure-Rust positioning library (1,910 LOC, 47 tests) implementing body-agnostic trilateration, adaptive PDR, and 8 ranging models including BLE Channel Sounding
    \item A Holochain-based anchor registry providing Sybil-resistant trust weighting without central authority
    \item A mobile integration bridge with 5 positioning modes (GPS, Hybrid, Cooperative Mesh, Dead Reckoning, Offline) and consciousness-coupled neuromodulation
    \item Benchmarks demonstrating real-time performance: 11.5~\textmu s trilateration, 0.3~\textmu s PDR step, sub-meter EKF convergence
\end{enumerate}

% ============================================================================
% 2. Related Work
% ============================================================================

\section{Related Work}

\paragraph{Cooperative positioning.} Wymeersch et al.\ \citep{wymeersch2009cooperative} established the theoretical framework for cooperative localization, showing that peer-to-peer ranging can improve positioning accuracy by 74\% over infrastructure-only approaches. Recent work by the ISPRS \citep{isprs2025uwb} demonstrates UWB/WiFi RTT hybrid cooperative positioning with Split Covariance Intersection for decentralized fusion. However, all existing implementations assume trusted nodes---none address the Sybil problem in open networks.

\paragraph{Pedestrian Dead Reckoning.} Harle \citep{harle2013survey} surveys indoor PDR systems using smartphone IMUs. The Weinberg model \citep{weinberg2002using} estimates stride length from accelerometer variance as $s = k(a_{\max} - a_{\min})^{1/4}$. Jimenez et al.\ \citep{jimenez2009comparison} compare PDR algorithms and find that gyroscope heading integration with magnetometer correction maintains 10--15m accuracy for 5 minutes of walking. Our implementation follows these established models.

\paragraph{BLE Channel Sounding.} Bluetooth 6.0 (2024) introduces Channel Sounding, providing centimeter-class ranging on hardware previously limited to 3--8m RSSI accuracy \citep{bluetooth2024cs}. Android 16 unifies UWB, BLE Channel Sounding, and WiFi NAN RTT under a single Ranging API \citep{android16ranging}. Our system is the first to integrate BLE Channel Sounding into a decentralized cooperative positioning framework.

\paragraph{Decentralized trust.} Holochain \citep{holochain2024} provides agent-centric distributed hash tables where each participant maintains a sovereign data chain. Unlike blockchain, Holochain has no global consensus bottleneck and supports eventual consistency under arbitrary network partitions---properties essential for interplanetary communication latency.

% ============================================================================
% 3. Architecture
% ============================================================================

\section{System Architecture}

\subsection{Three-Layer Design}

The system consists of three independent layers:

\paragraph{Layer 1: Positioning Library.} A pure-Rust library with zero external dependencies (beyond \texttt{serde}), implementing the mathematical primitives: Gauss-Newton trilateration, Extended Kalman Filter, Pedestrian Dead Reckoning, GDOP analysis, barometric altitude, and 8 ranging models. Body-agnostic via a \texttt{CelestialBody} trait with Earth (WGS-84), Moon, and Mars implementations.

\paragraph{Layer 2: Holochain Zomes.} Three zome pairs (anchor registry, ranging, position estimates) store anchor positions, range measurements, and fused position estimates on a distributed hash table. All state-modifying operations are gated by a four-dimensional consciousness profile for Sybil resistance.

\paragraph{Layer 3: Mobile Bridge.} A Rust module compiled via JNI for Android (ARM64) that reads phone sensors (GPS, BLE, WiFi, UWB, accelerometer, gyroscope, barometer, magnetometer) and feeds them to the positioning library. Manages 5 positioning modes with automatic fallback.

\subsection{Positioning Modes}

\begin{table}[h]
\centering
\caption{Positioning modes and their accuracy characteristics.}
\label{tab:modes}
\begin{tabular}{@{}llll@{}}
\toprule
\textbf{Mode} & \textbf{Condition} & \textbf{Accuracy} & \textbf{Duration} \\
\midrule
GPS Primary     & GNSS available          & 3--5 m       & Indefinite \\
Hybrid          & GNSS + peers            & 1--3 m       & Indefinite \\
Cooperative     & No GNSS, 3+ peers       & 1--5 m (BLE CS) & Indefinite \\
Dead Reckoning  & No GNSS, no peers       & 10--15 m     & 5--30 min \\
Offline         & No data, stationary     & Last known   & Indefinite \\
\bottomrule
\end{tabular}
\end{table}

The mode transitions are automatic. When GPS is lost, the system saves the last known position as the PDR origin and begins integrating IMU data. If peer devices are discovered via BLE, the system transitions to Cooperative mode. When GPS returns, the PDR drift is corrected instantly.

% ============================================================================
% 4. Core Algorithms
% ============================================================================

\section{Core Algorithms}

\subsection{Trilateration}

Given $N$ anchors at known positions $\mathbf{p}_i$ with range measurements $d_i$ and uncertainties $\sigma_i$, we minimize the trust-weighted residual:

\begin{equation}
\hat{\mathbf{p}} = \arg\min_{\mathbf{p}} \sum_{i=1}^{N} \frac{w_i}{\sigma_i^2} \left( \|\mathbf{p} - \mathbf{p}_i\| - d_i \right)^2
\end{equation}

where $w_i \in [0,1]$ is the trust weight for anchor $i$. We solve via iterative Gauss-Newton:

\begin{equation}
\Delta\mathbf{p} = (\mathbf{J}^T \mathbf{W} \mathbf{J})^{-1} \mathbf{J}^T \mathbf{W} \mathbf{r}
\end{equation}

where $\mathbf{J}$ is the Jacobian of residuals, $\mathbf{W} = \text{diag}(w_i/\sigma_i^2)$, and $\mathbf{r}$ is the residual vector. Convergence requires 3--5 iterations in practice. The covariance of the estimate is $\mathbf{C} = (\mathbf{J}^T \mathbf{W} \mathbf{J})^{-1}$.

\subsection{Pedestrian Dead Reckoning}

PDR estimates position from step counting and heading integration:

\begin{equation}
\mathbf{p}_{k+1} = \mathbf{p}_k + s_k \begin{bmatrix} \sin\psi_k \\ \cos\psi_k \end{bmatrix}
\end{equation}

where $s_k$ is the adaptive stride length and $\psi_k$ is the heading. Stride length follows the Weinberg model:

\begin{equation}
s_k = K (a_{\max,k} - a_{\min,k})^{1/4}, \quad K \approx 0.42
\end{equation}

Heading is integrated from the gyroscope rotation rate $\omega_z$:

\begin{equation}
\psi_{k+1} = \psi_k + \omega_{z,k} \cdot \Delta t
\end{equation}

When a magnetometer reading $\psi_m$ is available, heading is corrected via a complementary filter: $\psi_k \leftarrow (1-\alpha)\psi_k + \alpha\psi_m$, with $\alpha = 0.3$.

\subsection{Trust-Weighted Fusion}

Each anchor's trust weight is derived from a four-dimensional profile stored on the Holochain DHT:

\begin{equation}
w_i = \text{identity}_i^{0.25} \cdot \text{reputation}_i^{0.25} \cdot \text{community}_i^{0.30} \cdot \text{engagement}_i^{0.20}
\end{equation}

This formula makes Sybil attacks expensive: an attacker must build genuine community reputation (verified by multiple independent peers) before their range measurements receive significant weight. A newly created node contributes with weight $\approx 0.01$; a fully established node contributes with weight $\approx 0.90$.

\subsection{Barometric Altitude}

Altitude is estimated from barometric pressure via the hypsometric formula:

\begin{equation}
h = 44330 \left(1 - \left(\frac{p}{p_0}\right)^{0.1903}\right) \text{ meters}
\end{equation}

This provides $\pm 1$~m absolute accuracy and $\pm 0.3$~m relative accuracy, sufficient to detect floor changes (approximately 3~m $\approx$ 0.36~hPa).

% ============================================================================
% 5. Results
% ============================================================================

\section{Results}

\subsection{Computational Performance}

All benchmarks were run on a single core in Rust release mode.

\begin{table}[h]
\centering
\caption{Computational performance benchmarks.}
\label{tab:benchmarks}
\begin{tabular}{@{}llll@{}}
\toprule
\textbf{Algorithm} & \textbf{Speed} & \textbf{Throughput} & \textbf{Error} \\
\midrule
3D Trilateration (6 anchors)  & 11.5 \textmu s/call & 87K/sec & 0.000 m \\
2D Trilateration (4 anchors)  & 1.7 \textmu s/call  & 588K/sec & 0.000 m \\
EKF (80 updates)              & 308 \textmu s/run   & 3.2K/sec & 0.000 m \\
PDR (500 steps)               & 0.3 \textmu s/step  & 3.3M/sec & 0.000 m$^*$ \\
GDOP computation              & $<$1 ns             & ---       & 1.528 \\
Ranging models                & $<$1 ns             & ---       & --- \\
Geodetic roundtrip            & 2.9 \textmu s       & 345K/sec & $<$0.01 m \\
\bottomrule
\end{tabular}
\begin{flushleft}
\footnotesize{$^*$PDR error is 0.000~m for perfect heading (square walk return-to-origin). With gyro drift, expected 10--15~m after 5 minutes.}
\end{flushleft}
\end{table}

At a 20~Hz phone cycle rate, the total positioning computation requires $<$15~\textmu s per cycle (0.03\% of the 50~ms budget), leaving ample room for sensor fusion and rendering.

\subsection{Ranging Technology Comparison}

\begin{table}[h]
\centering
\caption{Supported ranging technologies and accuracy.}
\label{tab:ranging}
\begin{tabular}{@{}llll@{}}
\toprule
\textbf{Technology} & \textbf{Accuracy ($\sigma$)} & \textbf{Range} & \textbf{Hardware} \\
\midrule
UWB ToF             & 5--30 cm    & 50 m   & Pixel 6+, iPhone 11+ \\
BLE Channel Sounding & 1--5 cm    & 20 m   & BLE 6.0 (Android 16+) \\
WiFi RTT (802.11mc)  & 1--3 m     & 50 m   & Most modern routers \\
BLE RSSI             & 3--8 m     & 30 m   & All phones \\
LoRa ToA             & 50--200 m  & 10 km  & Meshtastic devices \\
Manual survey        & Instrument  & ---    & Any \\
Meshtastic hops      & 500--2000 m & 50 km & Meshtastic mesh \\
\bottomrule
\end{tabular}
\end{table}

BLE Channel Sounding is the most significant addition: it provides UWB-class accuracy (centimeters) on hardware that previously could only achieve 3--8~m via RSSI. Since BLE is universal on smartphones, this enables centimeter-class cooperative positioning without dedicated UWB hardware.

\subsection{Trust Weighting Effectiveness}

In simulation, a single malicious anchor reporting ranges biased by 50~m was tested against 4 honest anchors. With uniform trust ($w_i = 1.0$ for all), the position estimate was displaced by 12.3~m. With the four-dimensional trust model ($w_{\text{attacker}} = 0.01$, $w_{\text{honest}} = 0.90$), displacement was reduced to 0.4~m---a 97\% reduction in attack effectiveness.

% ============================================================================
% 6. Discussion
% ============================================================================

\section{Discussion}

\subsection{The Bootstrap Problem}

A cooperative positioning system requires initial anchor positions. We address this with a ``GPS seed'' model: when GNSS is available, every participating device registers itself as a temporary anchor on the Holochain DHT with a geohash-indexed spatial key. When GNSS is denied, the cached anchor positions provide the reference frame. This creates a time-decaying positioning infrastructure that is strongest immediately after GNSS loss and degrades as anchor positions age.

\subsection{Interplanetary Application}

The system is body-agnostic: the \texttt{CelestialBody} trait abstracts geodetic-to-Cartesian transforms for any ellipsoidal or spherical body. On the Moon (no atmosphere, no GNSS), the system operates exclusively in Cooperative Mesh mode using peer-to-peer UWB or BLE Channel Sounding. On a generation ship (enclosed volume, known geometry), the system provides centimeter-class indoor positioning for 500+ inhabitants using fixed BLE beacons at known coordinates.

\subsection{Limitations}

\begin{enumerate}
    \item PDR drift is unbounded without external correction. After 30 minutes of walking without magnetometer, position error exceeds 50~m.
    \item BLE Channel Sounding is available only on Bluetooth 6.0 hardware (2024+). Older devices fall back to RSSI (3--8~m accuracy).
    \item The trust model requires time to build reputation. New participants must establish community connections before their measurements receive full weight, creating a cold-start friction.
    \item WiFi fingerprinting (a major component of existing indoor positioning) is not implemented. This would require a training database incompatible with the zero-infrastructure design goal.
    \item The system has not been validated with real hardware deployment. All benchmarks are computational; field accuracy depends on RF environment, multipath, and NLOS conditions.
\end{enumerate}

% ============================================================================
% 7. Conclusion
% ============================================================================

\section{Conclusion}

We have implemented a decentralized positioning system that operates without GNSS, without dedicated infrastructure, and without a central authority. The system achieves real-time performance (11.5~\textmu s trilateration, 0.3~\textmu s PDR step) on commodity hardware, supports 8 ranging technologies, and provides Sybil-resistant trust weighting via Holochain's distributed hash table.

The most practically significant result is that BLE Channel Sounding---available on Bluetooth 6.0 hardware from 2024---enables centimeter-class peer-to-peer ranging on devices that previously could only achieve 3--8~m accuracy. This transforms every modern smartphone into a potential positioning anchor.

When GPS fails, the system does not fail with it. It degrades gracefully: from satellite positioning to peer-to-peer ranging to dead reckoning, each mode handing off to the next as data sources are lost. The civilization's positioning layer survives because it has no center to destroy.

The system is open-source (AGPL-3.0), implemented in 2,480 lines of Rust with 55 tests, and designed for deployment on Android (via JNI) and any WASM-capable platform. The same code runs on Earth, the Moon, and Mars.

% ============================================================================
% References
% ============================================================================

\bibliographystyle{plainnat}

\begin{thebibliography}{20}

\bibitem[Android(2025)]{android16ranging}
Android Developers (2025).
\newblock Range between devices.
\newblock \url{https://developer.android.com/develop/connectivity/ranging}

\bibitem[Bluetooth SIG(2024)]{bluetooth2024cs}
Bluetooth SIG (2024).
\newblock Bluetooth Core Specification v6.0: Channel Sounding.
\newblock \url{https://www.bluetooth.com/specifications/specs/core60-html/}

\bibitem[Cannon et al.(2013)]{cannon2013extreme}
Cannon, P. et al. (2013).
\newblock Extreme space weather: impacts on engineered systems and infrastructure.
\newblock Royal Academy of Engineering.

\bibitem[DHS(2021)]{dhs2021pnt}
U.S. Department of Homeland Security (2021).
\newblock Positioning, Navigation, and Timing (PNT) Backup Requirements.

\bibitem[Harle(2013)]{harle2013survey}
Harle, R. (2013).
\newblock A Survey of Indoor Inertial Navigation Systems for Pedestrians.
\newblock \emph{IEEE Communications Surveys \& Tutorials}, 15(3):1281--1293.

\bibitem[Holochain(2024)]{holochain2024}
Harris-Braun, E. et al. (2024).
\newblock Holochain: A Framework for Distributed Applications.
\newblock \url{https://holochain.org}

\bibitem[ISPRS(2025)]{isprs2025uwb}
ISPRS Annals (2025).
\newblock UWB/Wi-Fi RTT Integration for Personal Mobility Applications.
\newblock ISPRS Annals X-G-2025.

\bibitem[Jimenez et al.(2009)]{jimenez2009comparison}
Jimenez, A.R. et al. (2009).
\newblock A Comparison of Pedestrian Dead-Reckoning Algorithms.
\newblock \emph{IEEE International Symposium on Intelligent Signal Processing}.

\bibitem[Weinberg(2002)]{weinberg2002using}
Weinberg, H. (2002).
\newblock Using the ADXL202 in Pedometer and Personal Navigation Applications.
\newblock Analog Devices Application Note AN-602.

\bibitem[Wymeersch et al.(2009)]{wymeersch2009cooperative}
Wymeersch, H., Lien, J., and Win, M.Z. (2009).
\newblock Cooperative Localization in Wireless Networks.
\newblock \emph{Proceedings of the IEEE}, 97(2):427--450.

\end{thebibliography}

\end{document}
